\exnumber{00}
\extitle{UI III - Ex 1 - Grids}
\exfiles{ex00/Figma file, named Name\_FirstName\_DeliverableName\_Date\_VersionNumber}
\exauthorize{Those named in the exercise}
\makeheaderfiles

\textbf{Context:}

Welcome to the 3rd quest of UI.

In this series of exercises, you'll be asked to open your eyes and collect graphic elements from various websites. You'll see grids across screens, multi state buttons, how to make a design accessible for all, breadcrumbs, toggle buttons, radio buttons, calendar, time picker with 2 main goals:

\begin{itemize}
\item Enrich your graphic general knowledge
\item Detect consistency, that helps users understand a digital product at the first glance.
\end{itemize}

\textbf{Instructions:}

Open these 5 images on Figma and draw lines and grids every time you find aligned elements.

\textbf{Deliverables:}

\begin{itemize}
\item A Figma file including the screens with lines on top to show grids. One page per screen or frame.
\item Upload the Figma file on Github, labeled as follows: "Name\_FirstName\_DeliverableName\_Date\_VersionNumber". For example, it can be labeled as "Doe\_John\_MidFiPrototype\_05242024\_V1".
\item Add a title within the file.
\item Pay attention to the global aspect of the file. It must be clear, simple and easy to read. You can get inspiration on canva.com but don't overload your design with too much details!
\end{itemize}

\textbf{Resources:}

\begin{itemize}
\item Video tutorial on Figma \url{https://www.youtube.com/watch?v=FTFaQWZBqQ8}
\item Beginning graphic design \url{https://edu.gcfglobal.org/en/beginning-graphic-design/}
\item Gareth graphic studio \url{https://www.youtube.com/c/GarethDavidStudio}
\item 1st screen example \url{https://www.figma.com/file/KOzd8yZuJvGLmZLuXjBgDg/UI-III-Ex-1.1} 
\end{itemize}

